\section{Reversible Embedding versus Erased XOR}
\label{sec:results-reversible}

Table~\ref{tab:reversible-erased} isolates the accounting point of the paper in its cleanest form. All three rows come from the same underlying CNOT laboratory at the same noise setting. What changes is not the substrate but the package presented to the layer. The micro row keeps the full \(12\)-state laboratory. The retained macro row packages those microstates into the full four logical CNOT macrostates \((a,b)\). The XOR-erased row keeps only the one-bit output lens and forgets the retained control bit. In Six Birds terms, the split happens at P5: retaining or discarding side information changes the carrier of the layer. P6 then records the consequences.

\begin{table}[t]
    \centering
    \caption{Three views of the same CNOT laboratory: the full microstate view, the retained four-state CNOT macro view, and the XOR-erased one-bit output view. The table reports closure defect, route mismatch, entropy-production rate, and channel-information quantities, showing how forgetting the retained control bit changes both closure and retained information. Rendered directly from \protect\path{results/final_claims/table_reversible_vs_erased.csv}.}
    \label{tab:reversible-erased}
    {\small
\setlength{\tabcolsep}{4pt}
\renewcommand{\arraystretch}{1.12}
\begin{tabularx}{\linewidth}{@{}>{\raggedright\arraybackslash}X r r r r@{}}
\toprule
\multicolumn{5}{@{}l}{\textit{Closure and audit metrics}} \\
\midrule
{View} & {States} & {Defect} & {RM} & {EPR} \\
\midrule
CNOT micro & 12 & 0 & 0 & \(2.72\times 10^{-17}\) \\
CNOT macro & 4 & \(3.47\times 10^{-18}\) & \(1.70\times 10^{-16}\) & 0 \\
XOR erased macro & 2 & 0.96 & 0.96 & 0 \\
\bottomrule
\end{tabularx}

\medskip

\begin{tabularx}{\linewidth}{@{}>{\raggedright\arraybackslash}X r r r r r r@{}}
\toprule
\multicolumn{7}{@{}l}{\textit{Channel-information metrics}} \\
\midrule
{View} & {$H_{\mathrm{in}}$} & {$H_{\mathrm{out}}$} & {\makecell{$H(\mathrm{out}\mid$\\$\mathrm{in})$}} & {\makecell{$I(\mathrm{in};$\\$\mathrm{out})$}} & {$\Delta H$} & {$U_{\text{loss}}$} \\
\midrule
CNOT micro & 2 & 2 & 0.141 & 1.859 & 0 & 0.141 \\
CNOT macro & 2 & 2 & 0.141 & 1.859 & 0 & 0.141 \\
XOR erased macro & 2 & 1 & 0.141 & 0.859 & 1 & 1.141 \\
\bottomrule
\end{tabularx}
}

\end{table}

The first comparison to notice is that the micro and retained macro views are effectively the same logical object at different levels of description. The micro view has closure defect \(0\), RM \(0\), and a numerically tiny entropy-production rate \(2.72\times 10^{-17}\). The retained macro view has closure defect \(3.47\times 10^{-18}\), RM \(1.70\times 10^{-16}\), and apparent EPR \(0\). More importantly, their channel information measures agree at the logical level: both have \(H_{\mathrm{in}}=2\), \(H_{\mathrm{out}}=2\), \(H(\mathrm{out}\mid\mathrm{in})=0.141\), \(I(\mathrm{in};\mathrm{out})=1.859\), and \(U_{\text{loss}}=0.141\). The retained macro row is therefore not an information-poor shadow of the micro description. It is the same reversible embedding expressed at the carrier level that actually closes.

The erased XOR row is qualitatively different. Once the retained control bit is forgotten, the closure defect jumps to \(0.96\) and the RM jumps to \(0.96\). The frozen claims bundle records these as deltas of \(0.96\) and \(0.96\) relative to the retained macro view. At the same time, the channel information profile changes from a two-bit reversible embedding to a many-to-one reported output: \(H_{\mathrm{out}}\) drops from \(2\) to \(1\), \(I(\mathrm{in};\mathrm{out})\) drops from \(1.859\) to \(0.859\), and \(U_{\text{loss}}\) rises from \(0.141\) to \(1.141\). The resulting ratio of unretained input information is \(8.07\). This is the main accounting result of the section.

Two features of Table~\ref{tab:reversible-erased} are especially instructive. First, the conditional entropy \(H(\mathrm{out}\mid\mathrm{in})\) remains \(0.141\) in all three rows. The local noise on the reported output bit has not changed. What changes is how much of the input structure the reported channel can still carry once the control bit is forgotten. This is why the correct comparison is not simply entropy drop but unretained input information. Second, the apparent EPR of the two macro views is \(0\) in both cases. The decisive audit signal here is therefore not a large coarse thermodynamic signature; it is the joint change in closure and retained information. Forgetting side information does not merely add a surcharge to the same logical object. It produces a different logical object.

That distinction is the core P5/P6 lesson. Packaging is not an innocent matter of notation. The retained macro row keeps the full reversible carrier and therefore supports a closed two-bit operator. The erased XOR row reports only a one-bit summary and thereby turns a reversible embedding into a many-to-one channel with large closure failure. What looks superficially like ``the same XOR gate viewed more coarsely'' is, at the level of the audited layer, a different object with different closure properties and a much larger loss ledger.

The point matters beyond this one example. Many discussions of physical logic treat reversible and irreversible descriptions as though the latter were simply the former with some bookkeeping omitted. Table~\ref{tab:reversible-erased} shows why that is too weak. Omitting the bookkeeping changes the closure status of the layer itself. The next section turns from this comparison of declared views to the inverse problem: whether stable packaged bits and a gate can be recovered directly from transition structure without declaring the relevant logical objects in advance (Section~\ref{sec:results-discovery}).
